\section{Equivalence of potentials}

Throughout, $S$ is a countable parameter set, $(E,\mathcal{E})$ a measurable space,
$\Cfg = E^{S}$, $\Fin$ the collection of finite subsets of $S$, and $\Tail_\Lambda$ the
$\sigma$-algebra generated by the coordinates outside $\Lambda$. For a potential
$\Phi = (\Phi_A)_{A \in \Fin}$, an inverse temperature $\beta \in \R$ and an a priori measure
$\lambda$ on $(E,\mathcal{E})$ we write, as in \S2.1,
\[
  H_\Lambda^\Phi = \sum_{A \cap \Lambda \neq \emptyset} \Phi_A, \qquad
  h_\Lambda^\Phi = e^{-\beta H_\Lambda^\Phi}, \qquad
  Z_\Lambda^\Phi = \lambda_\Lambda h_\Lambda^\Phi, \qquad
  \rho_\Lambda^\Phi = h_\Lambda^\Phi / Z_\Lambda^\Phi ,
\]
and $\gamma^\Phi = \rho^\Phi \lambda$ for the associated $\lambda$-specification; $\beta = 1$ is
Georgii's normalization. Any factor of $h_\Lambda^\Phi$ that does not depend on the configuration
inside $\Lambda$ reappears in $Z_\Lambda^\Phi$ and cancels in $\rho_\Lambda^\Phi$;
Definition~(2.33) isolates the pairs of potentials differing by such a factor.

\begin{definition}[Equivalence of potentials, Georgii (2.33)]
    \label{def:equiv-potentials}
    \uses{def:cylinder-event}
    \lean{Potential.IsEquivalent}
    \leanok

    Two potentials $\Phi$, $\Psi$ are \textbf{equivalent}, written $\Phi \sim \Psi$, if for every
    $\Lambda \in \Fin$ the Hamiltonian $H_\Lambda^{\Phi - \Psi}$ is $\Tail_\Lambda$-measurable.
    The condition is stated for arbitrary families $\Phi, \Psi : \Fin \to \Cfg \to \R$ and carries
    Georgii's meaning under the standing summability hypothesis (2.2)(ii) on $\Phi$ and $\Psi$.
\end{definition}

\begin{lemma}[$\sim$ is an equivalence relation]
    \label{lem:equiv-equivalence-relation}
    \uses{def:equiv-potentials}
    \lean{Potential.IsEquivalent.refl, Potential.IsEquivalent.symm, Potential.IsEquivalent.trans}
    \leanok

    On summable potentials, $\sim$ is reflexive, symmetric and transitive.
\end{lemma}
\begin{proof}
    \leanok

    For summable $\Phi, \Psi, \Theta$ one has $H_\Lambda^{\Phi - \Psi} = H_\Lambda^{\Phi} -
    H_\Lambda^{\Psi}$ termwise, hence $H_\Lambda^{\Phi-\Phi} = 0$, $H_\Lambda^{\Psi - \Phi} =
    -H_\Lambda^{\Phi - \Psi}$ and $H_\Lambda^{\Phi - \Theta} = H_\Lambda^{\Phi-\Psi} +
    H_\Lambda^{\Psi-\Theta}$. Constants, negatives and sums of $\Tail_\Lambda$-measurable
    functions are $\Tail_\Lambda$-measurable.
\end{proof}

\begin{remark}[Adding constants to a potential]
    \label{rmk:equiv-add-constants}
    \uses{def:equiv-potentials}

    Let $c = (c_A)_{A \in \Fin}$ be reals with $\sum_{A \ni i} |c_A| < \infty$ for all $i \in S$
    and $\Psi_A = \Phi_A + c_A$. Then $H_\Lambda^{\Psi} - H_\Lambda^{\Phi} =
    \sum_{A \cap \Lambda \neq \emptyset} c_A$ is a finite constant, so $\Psi$ is
    $\lambda$-admissible whenever $\Phi$ is, and $\Phi \sim \Psi$. For $c_A = -\Phi_A(\eta_0)$,
    the case of Lemma~\ref{lem:equiv-centre}, summability of $c$ follows from summability of
    $\Phi$.
\end{remark}

\begin{lemma}[The Boltzmann factor of a difference splits off]
    \label{lem:equiv-boltzmann-mul}
    \lean{Potential.boltzmannFactor_eq_mul_sub}
    \leanok

    For summable $\Phi, \Psi$, every $\beta \in \R$, $\Lambda \in \Fin$ and $\eta \in \Cfg$,
    \[ h_\Lambda^{\Phi}(\eta) = h_\Lambda^{\Phi - \Psi}(\eta) \, h_\Lambda^{\Psi}(\eta). \]
\end{lemma}
\begin{proof}
    \leanok

    $H_\Lambda^{\Phi - \Psi} = H_\Lambda^\Phi - H_\Lambda^\Psi$ for summable potentials, and
    $e^{-\beta(a-b)}e^{-\beta b} = e^{-\beta a}$.
\end{proof}

\begin{lemma}[The difference factor is measurable outside the volume]
    \label{lem:equiv-boltzmann-measurable}
    \uses{def:equiv-potentials, def:cylinder-event}
    \lean{Potential.measurable_boltzmannFactor_sub_of_isEquivalent}
    \leanok

    If $\Phi \sim \Psi$ then $h_\Lambda^{\Phi - \Psi}$ is $\Tail_\Lambda$-measurable, for every
    $\beta \in \R$ and every $\Lambda \in \Fin$.
\end{lemma}
\begin{proof}
    \leanok

    It is the composition of the $\Tail_\Lambda$-measurable $H_\Lambda^{\Phi-\Psi}$ with the
    continuous map $x \mapsto e^{-\beta x}$, read in $[0,\infty]$.
\end{proof}

\begin{theorem}[Georgii (2.34), (i) $\Rightarrow$ (ii)]
    \label{thm:equiv-density-eq}
    \uses{def:equiv-potentials, lem:equiv-boltzmann-mul, lem:equiv-boltzmann-measurable, def:markov-spec}
    \lean{Potential.sigmaFinitePremodifierNorm_eq_of_isEquivalent}
    \leanok

    Let $\lambda \in \Meas(E, \mathcal{E})$ be $\sigma$-finite and let $\Phi, \Psi$ be summable
    potentials with $\Phi \sim \Psi$. Then $\rho^\Phi = \rho^\Psi$, for every $\beta \in \R$ and
    with no $\lambda$-admissibility assumption on $\Phi$ or $\Psi$: the identity is a
    cancellation in $[0,\infty]$, valid also where the partition functions degenerate to $0$ or
    $\infty$.
\end{theorem}
\begin{proof}
    \uses{lem:proper-ker-lintegral}
    \leanok

    Write $g = h_\Lambda^{\Phi-\Psi}$, which by Lemma~\ref{lem:equiv-boltzmann-measurable} is
    $\Tail_\Lambda$-measurable and takes values in $(0,\infty)$. By
    Lemma~\ref{lem:equiv-boltzmann-mul}, $h_\Lambda^\Phi = g \, h_\Lambda^\Psi$, and properness of
    the reference kernel $\lambda_\Lambda$ (Lemma~\ref{lem:proper-ker-lintegral}) pulls the factor
    $g$ out of the integral: $Z_\Lambda^\Phi(\eta) = g(\eta) Z_\Lambda^\Psi(\eta)$. Dividing,
    $g(\eta)$ cancels in $\rho_\Lambda^\Phi(\eta) = h_\Lambda^\Phi(\eta)/Z_\Lambda^\Phi(\eta)$.
\end{proof}

\begin{theorem}[Georgii (2.34), (i) $\Rightarrow$ (iii)]
    \label{thm:equiv-spec-eq}
    \uses{thm:equiv-density-eq, def:equiv-potentials, def:spec, def:premodif}
    \lean{Potential.lambdaSpecification_eq_of_isEquivalent}
    \leanok

    Let $\lambda \in \Meas(E,\mathcal{E})$ be non-zero and $\sigma$-finite, let $\beta \in \R$,
    and let $\Phi \sim \Psi$ be summable potentials such that $h^\Phi$ and $h^\Psi$ are
    pre-modifications which are $\lambda$-admissible, $0 < Z_\Lambda < \infty$ everywhere. Then
    $\gamma^\Phi = \gamma^\Psi$, where $\gamma^\Phi$ is the $\lambda$-specification of
    Definition~(1.27) attached to $h^\Phi$ through Remark~(1.32).
\end{theorem}
\begin{proof}
    \leanok

    The two specifications have the same reference kernels and, by
    Theorem~\ref{thm:equiv-density-eq}, the same densities $\rho^\Phi = \rho^\Psi$; two
    specifications with equal kernels at every $\Lambda$ and $\eta$ are equal.
\end{proof}

\begin{corollary}[Equivalent potentials have the same Gibbs measures]
    \label{cor:equiv-gibbs-eq}
    \uses{thm:equiv-spec-eq, def:gibbs-meas}

    If $\Phi \sim \Psi$ are $\lambda$-admissible then $\Gibbs{\Phi} = \Gibbs{\Psi}$; in
    particular (iii) $\Rightarrow$ (iv) whenever $\Gibbs{\Phi} \cup \Gibbs{\Psi} \neq \emptyset$.
\end{corollary}

\begin{lemma}[Georgii (2.34), (ii) $\Rightarrow$ (i): the logarithmic identity]
    \label{lem:equiv-log-partition}
    \lean{Potential.hamiltonian_sub_eq_log_sigmaFiniteLambdaZ}
    \leanok

    Let $\Phi, \Psi$ be summable potentials, $\lambda$-admissible at $\beta = 1$ in the sense of
    Georgii's (2.7), i.e.\ $0 < Z_\Lambda(\eta) < \infty$ for all $\Lambda \in \Fin$ and
    $\eta \in \Cfg$, and suppose $\rho^\Phi = \rho^\Psi$. Then
    \[ H_\Lambda^{\Phi}(\eta) - H_\Lambda^{\Psi}(\eta)
        = \log Z_\Lambda^{\Psi}(\eta) - \log Z_\Lambda^{\Phi}(\eta). \]
\end{lemma}
\begin{proof}
    \leanok

    Admissibility makes $Z_\Lambda^\Phi(\eta), Z_\Lambda^\Psi(\eta)$ finite and strictly
    positive, so $\rho_\Lambda^\Phi(\eta) = \rho_\Lambda^\Psi(\eta)$ may be read in $\R$ as
    $e^{-H_\Lambda^\Phi(\eta)}/Z_\Lambda^\Phi(\eta) =
    e^{-H_\Lambda^\Psi(\eta)}/Z_\Lambda^\Psi(\eta)$. Taking logarithms and rearranging gives the
    identity.
\end{proof}

\begin{theorem}[Georgii (2.34), (ii) $\Rightarrow$ (i)]
    \label{thm:equiv-of-density-eq}
    \uses{lem:equiv-log-partition, def:equiv-potentials}
    \lean{Potential.dependsOn_hamiltonian_sub_of_sigmaFinitePremodifierNorm_eq}
    \leanok

    Let $\Phi, \Psi$ be summable potentials, $\lambda$-admissible at $\beta = 1$, with
    $\rho^\Phi = \rho^\Psi$. Then for every $\Lambda \in \Fin$ the function
    $H_\Lambda^{\Phi - \Psi}$ takes the same value at any two configurations agreeing off
    $\Lambda$; being measurable, it is therefore $\Tail_\Lambda$-measurable, and
    $\Phi \sim \Psi$.
\end{theorem}
\begin{proof}
    \leanok

    By Lemma~\ref{lem:equiv-log-partition}, $H_\Lambda^{\Phi-\Psi} = \log(Z_\Lambda^\Psi /
    Z_\Lambda^\Phi)$, and the partition functions $Z_\Lambda^\Phi, Z_\Lambda^\Psi$ are
    $\Tail_\Lambda$-measurable, being integrals over the coordinates in $\Lambda$ against a
    reference kernel that only reads the configuration off $\Lambda$. Hence the right-hand side
    is unchanged when the configuration inside $\Lambda$ is altered.
\end{proof}

\begin{theorem}[Georgii (2.34)]
    \label{thm:equiv-tfae}
    \uses{thm:equiv-density-eq, thm:equiv-spec-eq, thm:equiv-of-density-eq, cor:equiv-gibbs-eq, def:gibbs-meas}

    Let $\lambda \in \Meas(E,\mathcal{E})$ and let $\Phi, \Psi$ be $\lambda$-admissible
    potentials. Consider
    \begin{enumerate}
        \item[(i)] $\Phi \sim \Psi$;
        \item[(ii)] $\rho^\Phi = \rho^\Psi$;
        \item[(iii)] $\gamma^\Phi = \gamma^\Psi$;
        \item[(iv)] $\Gibbs{\Phi} \cap \Gibbs{\Psi} \neq \emptyset$.
    \end{enumerate}
    Then (i) $\Leftrightarrow$ (ii) $\Rightarrow$ (iii), by
    Theorems~\ref{thm:equiv-density-eq}, \ref{thm:equiv-of-density-eq} and
    \ref{thm:equiv-spec-eq}. If moreover $E$ carries a second countable topology with
    $\mathcal{E}$ its Borel $\sigma$-algebra, $\lambda$ is everywhere dense, and each
    $\Phi_A - \Psi_A$ is continuous for the product topology on $\Cfg$, then (iii) $\Rightarrow$
    (ii). If in addition $\Phi - \Psi$ is uniformly convergent and
    $\Gibbs{\Phi} \cup \Gibbs{\Psi} \neq \emptyset$, then (i) $\Leftrightarrow$ (iv).
\end{theorem}

\begin{lemma}[Georgii (2.36) for $\alpha = \delta_a$]
    \label{lem:equiv-vacuum-hamiltonian}
    \lean{Potential.hamiltonian_vacuum}
    \leanok

    Let $a \in E$ and let $\Theta$ be a summable gas potential with vacuum state $a$, i.e.\
    $\Theta_A(\omega) = 0$ whenever $\omega_i = a$ for some $i \in A$. Then for all
    $\Lambda \in \Fin$ and $\omega \in \Cfg$,
    \[ H_\Lambda^{\Theta}(\omega_\Lambda a_{S \setminus \Lambda})
        = \sum_{\emptyset \neq A \subseteq \Lambda} \Theta_A(\omega), \]
    which is Georgii's $\alpha_{S \setminus \Lambda} H_\Lambda^{\Theta}$ for the Dirac measure
    $\alpha = \delta_a$.
\end{lemma}
\begin{proof}
    \leanok

    A term $\Theta_A$ with $A \cap \Lambda \neq \emptyset$ and $A \not\subseteq \Lambda$ vanishes
    at $\omega_\Lambda a_{S\setminus\Lambda}$, because that configuration carries the vacuum state
    at any $i \in A \setminus \Lambda$; a term with $\emptyset \neq A \subseteq \Lambda$ is
    $\mathcal{F}_A$-measurable and so equals $\Theta_A(\omega)$. Hence the partial sums over
    $\Delta \supseteq \Lambda$ are eventually constant, equal to the right-hand side.
\end{proof}

\begin{theorem}[Georgii (2.35)(a), Dirac branch]
    \label{thm:equiv-gas-zero}
    \uses{lem:equiv-vacuum-hamiltonian, def:equiv-potentials}
    \lean{Potential.eq_zero_of_isGasPotential}
    \leanok

    Let $a \in E$ and let $\Theta$ be a summable gas potential with vacuum state $a$ --- that is,
    $\Theta$ is normalized by $\alpha = \delta_a$ --- such that $H_\Lambda^{\Theta}$ is
    $\Tail_\Lambda$-measurable for every $\Lambda \in \Fin$, i.e.\ $\Theta \sim 0$. Then
    $\Theta_A = 0$ for every $A \in \Fin$, i.e.\ every nonempty finite $A$; nothing is claimed at
    $A = \emptyset$, whose value enters no Hamiltonian. Georgii states (2.35)(a) for a general
    normalizing $\alpha \in \Prob(E,\mathcal{E})$, under the alternative hypothesis that
    $\Phi - \Psi$ is uniformly convergent.
\end{theorem}
\begin{proof}
    \leanok

    By $\Tail_\Lambda$-measurability, $H_\Lambda^\Theta(\omega_\Lambda a_{S\setminus\Lambda}) =
    H_\Lambda^\Theta(a_S) = 0$, the last equality because every term of the Hamiltonian vanishes
    at the constant vacuum configuration. So by Lemma~\ref{lem:equiv-vacuum-hamiltonian}
    $\sum_{\emptyset \neq A \subseteq B} \Theta_A(\omega) = 0$ for every $B \in \Fin$; induction
    on $|B|$ (the inclusion--exclusion of Georgii's proof) gives $\Theta_B(\omega) = 0$.
\end{proof}

\begin{theorem}[Uniqueness of the gas representative]
    \label{thm:equiv-gas-unique}
    \uses{thm:equiv-gas-zero, thm:equiv-of-density-eq}
    \lean{Potential.eq_of_isGasPotential_of_sigmaFinitePremodifierNorm_eq}
    \leanok

    Let $a \in E$ and let $\Phi, \Psi$ be summable gas potentials with the same vacuum state $a$,
    $\lambda$-admissible at $\beta = 1$, with $\rho^\Phi = \rho^\Psi$ (equivalently, by
    Theorems~\ref{thm:equiv-of-density-eq} and \ref{thm:equiv-density-eq}, $\Phi \sim \Psi$).
    Then $\Phi_A = \Psi_A$ for every $A \in \Fin$.
\end{theorem}
\begin{proof}
    \leanok

    $\Phi - \Psi$ is again a summable gas potential with vacuum state $a$ (Georgii (2.29)(1)). By
    Theorem~\ref{thm:equiv-of-density-eq} its Hamiltonians depend only on the configuration
    outside the volume, so Theorem~\ref{thm:equiv-gas-zero} applies and gives
    $(\Phi - \Psi)_A = 0$ on nonempty $A$.
\end{proof}

\begin{theorem}[Georgii (2.35)(b)]
    \label{thm:equiv-gas-representative}
    \uses{thm:equiv-gas-unique, thm:equiv-tfae, thm:grep-2-30}

    Let $\Phi$ be such that for some $\lambda \in \Meas(E,\mathcal{E})$ the $\lambda$-modification
    $\rho^\Phi$ is quasilocal. Then for each $a \in E$ there is exactly one gas potential $\Psi$
    with vacuum state $a$ such that $\Phi \sim \Psi$; it is Georgii's M\"obius expression
    $\Psi_A = -\sum_{C \subseteq A} (-1)^{|A \setminus C|} \alpha_{S \setminus C} \log \rho_A^\Phi$
    at $\alpha = \delta_a$. If $\Phi$ is moreover uniformly convergent with bounded Hamiltonians,
    then for each $\alpha \in \Prob(E,\mathcal{E})$ there is exactly one uniformly convergent
    $\alpha$-normalized $\Psi$ with $\Phi \sim \Psi$.

    Existence in the first assertion is Theorem~\ref{thm:grep-2-30} applied to $\rho^\Phi$;
    uniqueness is Theorem~\ref{thm:equiv-gas-unique}.
\end{theorem}

\begin{remark}[(2.37), and Example (2.38): spin potentials versus lattice gas potentials]
    \label{rmk:equiv-normalized-formula}
    \uses{thm:equiv-gas-representative}

    The $\alpha$-normalized representative is given in closed form by Georgii's (2.37),
    \[ \Psi_A = \sum_{B \supseteq A} \sum_{C \subseteq A} (-1)^{|A \setminus C|}
        \alpha_{S \setminus C} \Phi_B , \]
    whence $\Psi$ is supported on subsets of the supports of $\Phi$, so a pair or
    nearest-neighbour potential stays one. Example~(2.38) applies this on $E = \{-1,+1\}$: the
    spin potential $\Phi_A = -J(A) \sigma^A$ is equivalent to the lattice gas potential
    $\Psi_A = K(A) \prod_{i \in A}(\sigma_i+1)/2$ with
    $K(A) = -2^{|A|} \sum_{B \supseteq A} (-1)^{|B \setminus A|} J(B)$, and the estimates
    $\|\Psi\|_i \leq 2 \sum_{B \ni i} |J(B)| 3^{|B|-1}$ and
    $\|\Phi\|_i \leq \tfrac12 \|\Psi\|_i$ of (2.39)--(2.40) show that $\Psi \in \Bspace$ implies
    $\Phi \in \Bspace$, but not conversely.
\end{remark}

\subsection*{Recentring, and the oscillation seminorm}

The normalization $\|\Phi_A\| = \osc(\Phi_A)/2$ performed at the start of the proof of
Theorem~(8.39) is isolated here as a statement about equivalence classes. Throughout,
$\osc(f) = \sup_{\zeta,\eta} |f(\zeta) - f(\eta)|$ and
$\|\Phi\|_i = \sum_{A \ni i} \sup_\eta |\Phi_A(\eta)|$ is Georgii's seminorm (2.12), finite at
every $i$ exactly when $\Phi \in \Bspace$.

\begin{definition}[Recentring a potential]
    \label{def:equiv-centre}
    \lean{Potential.centre}
    \leanok

    For $\eta_0 \in \Cfg$, the potential $\Phi$ \textbf{recentred at} $\eta_0$ is
    $(\Phi^{\eta_0})_A = \Phi_A - \Phi_A(\eta_0)$. It is again a potential in the sense of
    (2.2)(i), and is summable whenever $\Phi$ is.
\end{definition}

\begin{lemma}[Recentring stays in the equivalence class]
    \label{lem:equiv-centre}
    \uses{def:equiv-centre, def:equiv-potentials, rmk:equiv-add-constants}
    \lean{Potential.hamiltonian_sub_centre, Potential.isEquivalent_centre}
    \leanok

    For summable $\Phi$ and any $\eta_0 \in \Cfg$ one has
    $H_\Lambda^{\Phi - \Phi^{\eta_0}}(\eta) = H_\Lambda^{\Phi}(\eta_0)$ for all $\eta$; in
    particular $\Phi \sim \Phi^{\eta_0}$.
\end{lemma}
\begin{proof}
    \leanok

    Termwise, $(\Phi - \Phi^{\eta_0})_A = \Phi_A(\eta_0)$ is constant in $\eta$, and the terms
    entering $H_\Lambda$ are exactly those with $A \cap \Lambda \neq \emptyset$; the resulting
    series is $H_\Lambda^{\Phi}(\eta_0)$. A constant function is $\Tail_\Lambda$-measurable.
\end{proof}

\begin{definition}[Total oscillation at a site]
    \label{def:equiv-osc-norm}
    \lean{Potential.oscNormAt}
    \leanok

    $\displaystyle \osc_i(\Phi) = \sum_{A \ni i} \osc(\Phi_A) \in [0,\infty]$, the total
    oscillation at $i$ of the interaction terms containing $i$. Unlike $\|\Phi\|_i$ it sees only
    the interaction, not the additive constants.
\end{definition}

\begin{lemma}[Elementary properties of the oscillation seminorm]
    \label{lem:equiv-osc-norm-basic}
    \uses{def:equiv-osc-norm, def:equiv-centre}
    \lean{Potential.oscNormAt_centre, Potential.normAt_centre_le, Potential.oscNormAt_le_two_mul_normAt, Potential.oscNormAt_ne_top}
    \leanok

    For every $\Phi$, $\eta_0 \in \Cfg$ and $i \in S$:
    \begin{enumerate}
        \item $\osc_i(\Phi^{\eta_0}) = \osc_i(\Phi)$;
        \item $\|\Phi^{\eta_0}\|_i \leq \osc_i(\Phi)$;
        \item $\osc_i(\Phi) \leq 2 \|\Phi\|_i$, so $\Phi \in \Bspace$ implies
              $\osc_i(\Phi) < \infty$.
    \end{enumerate}
\end{lemma}
\begin{proof}
    \leanok

    (1) $\osc$ is invariant under adding a constant. (2) $|\Phi_A(\eta) - \Phi_A(\eta_0)| \leq
    \osc(\Phi_A)$ termwise. (3) $|\Phi_A(\zeta) - \Phi_A(\eta)| \leq \sup |\Phi_A| + \sup
    |\Phi_A|$ termwise.
\end{proof}

\begin{theorem}[Recentring criterion]
    \label{thm:equiv-centre-abs-summable}
    \uses{lem:equiv-osc-norm-basic, def:equiv-osc-norm, def:equiv-centre, lem:equiv-centre}
    \lean{Potential.isAbsolutelySummable_centre_iff}
    \leanok

    Assume $E \neq \emptyset$. Some recentring $\Phi^{\eta_0}$ of $\Phi$ lies in $\Bspace$ if and
    only if $\osc_i(\Phi) < \infty$ for every $i \in S$. Since $\Phi \sim \Phi^{\eta_0}$
    (Lemma~\ref{lem:equiv-centre}), finiteness of the total oscillation at every site is
    sufficient for the equivalence class of $\Phi$ to meet $\Bspace$; it is not necessary
    (Example~\ref{ex:equiv-osc-not-necessary}).
\end{theorem}
\begin{proof}
    \leanok

    ($\Leftarrow$) Recentre at any constant configuration and use
    Lemma~\ref{lem:equiv-osc-norm-basic}(2). ($\Rightarrow$) If $\Phi^{\eta_0} \in \Bspace$ then
    $\osc_i(\Phi) = \osc_i(\Phi^{\eta_0}) \leq 2\|\Phi^{\eta_0}\|_i < \infty$
    by Lemma~\ref{lem:equiv-osc-norm-basic}(1) and (3).
\end{proof}

\begin{example}[Finiteness of the oscillation is not necessary for the class to meet $\Bspace$]
    \label{ex:equiv-osc-not-necessary}
    \uses{thm:equiv-centre-abs-summable, def:equiv-potentials}

    Take $S \supseteq \{0,1,2\}$, an unbounded measurable $f : E \to \R$, and the potential with
    $\Phi_{\{0,1\}}(\eta) = f(\eta_0)$, $\Phi_{\{0,2\}}(\eta) = -f(\eta_0)$ and all other terms
    zero. Every Hamiltonian $H_\Lambda^{\Phi}$ is $0$ or a single $\pm f(\eta_0)$ cancelling
    against its partner, so $\Phi \sim 0 \in \Bspace$, while $\osc_0(\Phi) = \infty$.
\end{example}
